% Part: many-valued-logic
% Chapter: three-valued-logics
% Section: lukasiewicz

\documentclass[../../../include/open-logic-section]{subfiles}

\begin{document}

\olfileid{mvl}{thr}{luk}

\olsection{لوکاشیویچ دی منطق}

کئی قدراں والی منطق دیاں پہلِیاں چھپیاں تے تفصیل نال تیار
کیتیاں تجویزاں وچوں اک پولینڈ دے فلسفی یان لوکاشیویچ دی 1921
دی تجویز اے۔ لوکاشیویچ نوں ارسطو دی سمندری لڑائی والے مسئلے
توں خیال آیا۔ لگدا اے کہ \emph{اجّ} جملہ «بھلکے سمندر وچ
اک لڑائی ہووے گی» نہ سچ اے نہ جھوٹ، کیوں جے اوہدی صداقتی قدر
اجے طے نہیں ہوئی۔ لوکاشیویچ نے ایہو جیہے جملے، جیہڑے مستقبل
دے اوہناں واقعات بارے نیں جو ہو وی سکدے نیں تے نہیں وی،
لئی تِیجی صداقتی قدر شامل کرن دی تجویز دِتّی۔
\begin{quote}
میں کسے تضاد توں بغیر منّ سکدا آں کہ اگلے سال کسے خاص ویلے،
مثلاً 21 دسمبر دوپہر نوں، وارسا وچ میرا ہونا ہُن نہ ہاں وچ
طے اے نہ ناں وچ۔ ایس لئی ممکن اے کہ میں دِتّے ویلے وارسا
وچ ہواں، پر ضروری نہیں۔ ایس مفروضے تحت قضیہ «میں اگلے سال
21 دسمبر دوپہر نوں وارسا وچ ہواں گا» ایس ویلے نہ سچ ہو سکدا
اے نہ جھوٹ۔ کیوں جے ایہ ہُن سچ ہووے، تاں میرا اگّوں وارسا
وچ ہونا ضروری ہووے گا، جیہڑا مفروضے دے اُلٹ اے۔ دوجے پاسے،
جے ایہ ہُن جھوٹ ہووے، تاں میرا اگّوں وارسا وچ ہونا ناممکن
ہووے گا، جیہڑا وی مفروضے دے اُلٹ اے۔ سو ایہ قضیہ ایس ویلے
نہ سچ اے نہ جھوٹ تے ایہدی تِیجی قدر ہونی چاہیدی اے، جیہڑی
«0» یعنی جھوٹ تے «1» یعنی سچ توں وکھری ہووے۔ ایہ قدر
اسیں «$\frac{1}{2}$» لکھ سکدے آں۔ ایہ «ممکن» نوں ظاہر
کردی اے تے تِیجی قدر وجوں «سچ» تے «جھوٹ» دے نال جُڑدی اے۔
\end{quote}
اسیں لوکاشیویچ دی تِیجی صداقتی قدر لئی $\Undef$ ورتاں گے۔
\footnote{لوکاشیویچ ایتھے «ممکن» نوں اوہناں معنیاں وچ ورتدا
اے جیہڑے اجّ عام نہیں، یعنی ممکن پر ضروری نہیں۔}

ایس سمجھ مطابق $\lnot$، $\land$ تے $\lor$ دے صداقتی
فنکشن طے کرنے سوکھے نیں۔ مستقبل بارے غیر طے شدہ جملے دی
نفی وی مستقبل بارے غیر طے شدہ جملہ اے، ایس لئی
$\tf{\lnot}(\Undef) = \Undef$۔ جے عطف دا اک جزو غیر طے
شدہ تے دوجا سچ اے، تاں عطف وی غیر طے شدہ اے۔ مستقبل والا
جزو جیویں وی نکلے، اوہدے مطابق عطف سچ وی نکل سکدا اے تے
جھوٹ وی۔ ایس لئی
\[
    \tf{\land}(\True, \Undef) =
\tf{\land}(\Undef, \True) =
\Undef.
\]
پر جے دوجا جزو جھوٹ اے، تاں عطف سچ نہیں نکل سکدا۔ ایس لئی
\[\tf{\land}(\False, \Undef) =
\tf{\land}(\False, \Undef) = \False.\]
باقی قدراں، جدوں دلیلاں دیاں قدراں طے ہون، یعنی $\True$
یا $\False$ ہون، کلاسیکی منطق والیاں نیں۔

شرطیے لئی گل ذرا گُنجھلدار اے۔ منّ لو $q$ مستقبل بارے غیر
طے شدہ جملہ اے۔ جے $p$ جھوٹ اے، تاں $p \lif q$ سچ ہووے گا،
بھاویں $q$ جیویں وی نکلے۔ سو $\tf{\lif}(\False, \Undef) = \True$
رکھنا چاہیدا اے۔ تے جے $p$ سچ اے، تاں $q \lif p$ سچ ہووے گا،
بھاویں $q$ جیویں وی نکلے۔ ایس لئی $\tf{\lif}(\Undef, \True) = \True$۔
جے $p$ سچ اے، تاں $p
\lif q$ سچ وی نکل سکدا اے تے جھوٹ وی، سو $\tf{\lif}(\True,
\Undef) = \Undef$۔ ایسے طرح جے $p$ جھوٹ اے، تاں $q
\lif p$ سچ وی نکل سکدا اے تے جھوٹ وی، سو
$\tf{\lif}(\Undef, \False) = \Undef$۔ ہُن اوہ صورت باقی اے
جدوں $p$ تے $q$ دوویں مستقبل بارے غیر طے شدہ نیں۔ مُڈھلے
خیال مطابق ایس صورت وچ $\Undef$ دینا چاہیدا اے۔ پر ایس
نال $!A \lif !A$ ہمیشہ صادق \emph{نہیں} رہے گا۔ لوکاشیویچ
نوں $!A \lor \lnot !A$ تے $\lnot(!A
\land \lnot !A)$ چھڈن وچ کوئی دِقّت نہیں سی، پر اوہ $!A \lif !A$
چھڈن لئی تیار نہیں سی۔ ایس لئی اوہنے مقرر کیتا کہ
$\tf{\lif}(\Undef, \Undef) =
\True$۔

\begin{defn}\ollabel{def:lukasiewicz}
تِن قدراں والی لوکاشیویچ دی منطق ایس میٹرکس نال متعین اے۔
\begin{enumerate}
  \item معیاری قضیاتی زبان $\Lang L_0$ جیہدے وچ
  $\lnot$، $\land$، $\lor$، $\lif$ نیں۔
  \item صداقتی قدراں دا سیٹ $V = \{\True, \Undef, \False\}$۔
  \item $\True$ اکّو مقرر قدر اے، یعنی $V^+ = \{\True\}$۔
  \item صداقتی فنکشن ہِٹھلے جدولاں توں ملدے نیں۔
  \begin{center}
    \begin{tabular}{c|c} 
      $\tf{\lnot}$ & \\ 
      \hline  
      $\True$ & $\False$ \\ 
      $\Undef$ & $\Undef$ \\
      $\False$ & $\True$ 
    \end{tabular}
    \quad
    \begin{tabular}{c|ccc} 
      $\tf{\land}[\LogLuk[3]]$ & $\True$ & $\Undef$ & $\False$ \\ 
      \hline 
      $\True$ & $\True$ & $\Undef$ & $\False$ \\ 
      $\Undef$ & $\Undef$ & $\Undef$ & $\False$\\ 
      $\False$ & $\False$ & $\False$ & $\False$ 
    \end{tabular}
    \\[2ex]
    \begin{tabular}{c|ccc} 
      $\tf{\lor}[\LogLuk[3]]$ & $\True$ & $\Undef$ & $\False$ \\ 
      \hline 
      $\True$ & $\True$ & $\True$ & $\True$ \\ 
      $\Undef$ & $\True$ & $\Undef$ & $\Undef$ \\
      $\False$ & $\True$ & $\Undef$ & $\False$ 
    \end{tabular}
    \quad
    \begin{tabular}{c|ccc} 
      $\tf{\lif}[\LogLuk[3]]$ & $\True$ & $\Undef$ & $\False$ \\ 
      \hline 
      $\True$ & $\True$ & $\Undef$ & $\False$ \\ 
      $\Undef$ & $\True$ & $\True$ & $\Undef$  \\ 
      $\False$ & $\True$ & $\True$ & $\True$ 
    \end{tabular}
  \end{center}
\end{enumerate}
\end{defn}

سوکھے ای دِسدا اے کہ کوئی وی !!{formula}~$!A$ جیہدے وچ
صرف $\lnot$، $\land$ تے $\lor$ ہوون، صداقتی قدر~$\Undef$
لیندا اے جے اوہدے سارے !!{propositional variable}s نوں~$\Undef$
دِتّی جاوے۔ سو مثلاً کلاسیکی ہمیشہ صادق فارمولے $p \lor \lnot p$
تے $\lnot(p
\land \lnot p)$، $\LogLuk[3]$ وچ ہمیشہ صادق نہیں، کیوں جے
$\pValue{v}(!A)
= \Undef$ جدوں $\pAssign v(p) = \Undef$۔

اوہناں !!{valuation}s وچ جتھے $\pAssign v(p) = \True$ یا
$\False$، $\pValue
v(!A)$ اپنی کلاسیکی صداقتی قدر دے برابر ہووے گا۔

\begin{prop}
  جے $\pAssign v(p) \in \{\True, \False\}$ ہر اوہ $p$ لئی
  جیہڑا~$!A$ وچ اے، تاں
  $\pValue v(!A)[\LogLuk[3]] = \pValue v(!A)[\LogCL]$۔
\end{prop}

\begin{prob}\label{mvl:thr:luk:prob:luk-iff} منّ لو اسیں
  $\pValue v(!A \liff !B) = \pValue v((!A \lif !B) \land (!B \lif
  !A))$ نوں~$\LogLuk[3]$ وچ تعریف وجوں ورتدے آں۔ $\liff$ دا
  صداقتی جدول کیہ ہووے گا؟
\end{prob}

کئی کلاسیکی ہمیشہ صادق فارمولے \LogLuk[3] وچ وی \emph{ہمیشہ
صادق} نیں، مثلاً $\lnot p \lif (p \lif q)$۔ کلاسیکی منطق وانگ
اسیں ایہ گل صداقتی جدولاں نال پرکھ سکدے آں۔
\begin{center}
\begin{tabular}{cc|cccccc}
  $p$ & $q$ & $\lnot$ & $p$ & $\lif$ & $\smash{(}p$ & $\lif$ & $q\smash{)}$ \\ \hline
  \True & \True & \False & \True & \True & \True & \True & \True \\
  \True & \Undef & \False & \True & \True & \True & \Undef & \Undef \\
  \True & \False & \False & \True & \True & \True & \False & \False \\
  \Undef & \True & \Undef & \Undef & \True & \Undef & \True & \True \\
  \Undef & \Undef & \Undef & \Undef & \True & \Undef & \True & \Undef \\
  \Undef & \False & \Undef & \Undef & \True & \Undef & \Undef & \False \\
  \False & \True & \True & \False & \True & \False & \True & \True \\
  \False & \Undef & \True & \False & \True & \False & \True & \Undef \\
  \False & \False & \True & \False & \True & \False & \True & \False \\  
\end{tabular}
\end{center}

\begin{prob}
  دکھاؤ کہ ہِٹھلے فارمولے \LogLuk[3] وچ ہمیشہ صادق نیں۔
  \begin{enumerate}
    \item $p \lif (q \lif p)$
    \item\label{mvl:thr:luk:prob:luk-taut-2} $\lnot(p \land q) \liff (\lnot p \lor \lnot q)$
    \item\label{mvl:thr:luk:prob:luk-taut-3} $\lnot(p \lor q) \liff (\lnot p \land \lnot q)$
  \end{enumerate}
  (\olref[mvl][thr][luk]{prob:luk-taut-2} تے
  \olref[mvl][thr][luk]{prob:luk-taut-3} وچ $!A \liff !B$
  نوں $(!A \lif !B) \land (!B \lif !A)$ دا مختصر روپ سمجھو،
  یا \olref[mvl][thr][luk]{prob:luk-iff} دا اپنا حل ویکھو۔)
\end{prob}

\begin{prob}
  دکھاؤ کہ ہِٹھلے کلاسیکی ہمیشہ صادق فارمولے~\LogLuk[3] وچ
  ہمیشہ صادق نہیں۔
  \begin{enumerate}
    \item $(\lnot p \land p) \lif q)$
    \item $((p \lif q) \lif p) \lif p$
    \item $(p \lif (p \lif q)) \lif (p \lif q)$
  \end{enumerate}
\end{prob}

سو شاید ایہ خیال آوے کہ بھاویں سارے کلاسیکی ہمیشہ صادق
فارمولے~$\LogLuk[3]$ وچ ہمیشہ صادق نہیں، فیر وی اوہناں دی قدر~$\True$ یا~$\Undef$ ہونی چاہیدی اے
ہر !!{valuation} اُتے۔ ایہ گل سچ نہیں۔ اک ردّی مثال ایہ اے۔
\[
  \lnot(p \lif \lnot p) \lor \lnot(\lnot p \lif p)
\]
ایہدی قدر $\False$ اے جے $p$ دی قدر~$\Undef$ اے۔

\begin{prob}
  ہِٹھلے تعلقاں وچوں کیہڑے لوکاشیویچ دی منطق وچ قائم نیں؟
  ہر اک لئی صداقتی جدول دیو۔
  \begin{enumerate}
    \item $p, p \lif q \Entails q$
    \item $\lnot\lnot p \Entails p$
    \item $p \land q \Entails p$
    \item $p \Entails p \land p$
    \item $p \Entails p \lor q$
  \end{enumerate}
\end{prob}

لوکاشیویچ نوں امید سی کہ اپنے تِن قدراں والے نظام اُتے امکان
دی منطق بنا سکے گا۔ ایس لئی اوہنے اک مقامی رابطہ $\Diamond
!A$، یعنی «$!A$ ممکن اے»، تے اوہدے نال ملدا $\Box !A$،
یعنی «$!A$ ضروری اے»، شامل کیتے۔
\begin{center}
  \begin{tabular}{c|c} 
    $\tf{\Diamond}$ & \\ 
    \hline  
    $\True$ & $\True$ \\ 
    $\Undef$ & $\True$ \\
    $\False$ & $\False$ 
  \end{tabular}
  \quad
  \begin{tabular}{c|c} 
    $\tf{\Box}$ & \\ 
    \hline  
    $\True$ & $\True$ \\ 
    $\Undef$ & $\False$ \\
    $\False$ & $\False$ 
  \end{tabular}
\end{center}
دوجے لفظاں وچ، $p$ اُدوں تے صرف اُدوں ممکن اے جدوں ایہ
پہلاں ای جھوٹ طے نہ ہووے۔ تے $p$ اُدوں تے صرف اُدوں ضروری
اے جدوں ایہ پہلاں ای سچ طے ہووے۔

\begin{prob}
  دکھاؤ کہ $\Box p \liff \lnot\Diamond \lnot p$ تے $\Diamond p \liff
  \lnot \Box \lnot p$، ~\LogLuk[3] وچ ہمیشہ صادق نیں جدوں
  اینوں $\Box$ تے~$\Diamond$ دے صداقتی جدولاں نال ودھایا جاوے۔
\end{prob}

پر ایس تجویز کیتی ہوئی امکانی منطق دیاں کمزوریاں چھیتی
سامھنے آ گئیاں۔ جیویں وی گل نکلے، $p \land \lnot p$ کدے سچ
نہیں نکل سکدا۔ سو بھاویں ایہ ہُن طے نہیں ہویا، تے ایس لئی
غیر طے شدہ اے، اینوں ناممکن منّنا چاہیدا اے۔ یعنی
$\lnot \Diamond(p
\land \lnot p)$ ہمیشہ صادق ہونا چاہیدا اے۔ پر جے $\pAssign v(p) =
\Undef$، تاں $\pValue v(\lnot \Diamond(p \land \lnot p)) =
\Undef$۔ لوکاشیویچ دی ایہ گل ٹھیک سی کہ امکان تے ضرورت
ورگے فرق رکھن لئی دو صداقتی قدراں کافی نہیں۔ پر تِیجی
صداقتی قدر شامل کرنا وی کافی نہیں۔

\end{document}
